\begin{afterword}
\summaryhead

The post distinguishes terminal values, wanted for themselves, from instrumental values, wanted for what they lead to. People keep the two apart when they plan: no one drives to a supermarket whose chocolate has been destroyed. But they confuse them in abstract discussion, partly because English uses one word, ``want'', for both. A first definition in English fails, since even saving one's sister would be given up if it caused the Earth's destruction.

The author therefore sets out an expected-utility model with four parts: outcomes, actions, a utility for each outcome, and a probability of each outcome given each action. The chosen action is one with the highest expected utility. Utility belongs to outcomes and expected utility to actions. So a philosopher who says that a mother who died for her son did it for ``the feeling of importance'' she attached to the decision has confused the two, a ``type error''. In the flat model instrumental values disappear. They return when causal structure is represented, as shortcuts that depend on beliefs of fact and should be dropped when they stop working.

Some moral disputes concern means and some concern terminal values. The author fears that the brain does not keep these apart, and concludes that we do not know our own values or where they came from.

\responsehead

The post states its model correctly. The expected-utility rule, the separation of utility from expected utility, and the worked example, with the sister's chances of 0.9 and 0.3, are right. The central distinction is old, as ends and means or as final and instrumental value, and the post does not claim otherwise. Its main use of the model, against the selfish philosopher, is sound as far as it goes. That the mother chose what she rated highest is true of any choice, and says nothing about what she valued. Philosophers make the same objection to psychological egoism.

Two formal steps need qualifying. First, a plan of several steps cannot always be one fixed sequence of acts. When something may be learned along the way, like the radio news in the post's own opening, each option must be a strategy that says what to do after each possible observation. The model still applies, over a larger space. Second, the reply to the philosopher catches a type error in one reading of the words. On another, the feeling of importance is part of the outcome: she preferred dying with it to living without it. The model allows that view, and what counts against it is evidence about motives, which the post does not give.

The claims about people rest on the author's observation. That abstract talk confuses people is something the author has ``noticed''. That the brain does not ``strongly type'' moral beliefs about means and ends is offered with ``I fear'', and supported by how such beliefs feel. The circumcision example, given as a dispute about terminal values, is mixed: the reasons for the practice that the World Health Organization reports include beliefs of fact as well as values its opponents reject. The division of moral disputes into disputes of fact and disputes of value was drawn by Charles Stevenson in 1944.

In short: a correct sketch of expected-utility theory, used soundly against one form of egoism, with a model of plans that holds only when nothing is learned on the way, and claims about moral psychology that rest on introspection.
\end{afterword}
